\paragraph{E.3r (Proved: a specified generic fourfold and a two-torsion
counterexample).}
We can remove the unspecified auxiliary prime and generality choices
from the variety in E.3h--k. Put
\[
 \mathcal M_2=\{\alpha\in\mathbb Z_{\ge0}^3:|\alpha|=2\},
 \qquad
 k=\overline{\mathbb Q(a_{j\epsilon\alpha},b_{j\epsilon\alpha})},
\]
where all $72$ elements $a_{j\epsilon\alpha},b_{j\epsilon\alpha}$,
with $1\le j\le3$, $0\le\epsilon\le1$, and
$\alpha\in\mathcal M_2$, are algebraically independent.
On $\mathbb P^1_k\times\mathbb P^2_k\times\mathbb P^2_k$,
with coordinates $[S:T]$, $[x_0:x_1:x_2]$, $[y_0:y_1:y_2]$, set
\[
 A_j=\sum_{\epsilon,\alpha}
      a_{j\epsilon\alpha}S^{1-\epsilon}T^\epsilon x^\alpha,
 \qquad
 B_j=\sum_{\epsilon,\alpha}
      b_{j\epsilon\alpha}S^{1-\epsilon}T^\epsilon y^\alpha.
\]
Define $W/k$ by the three $2$-by-$2$ minors of
\[
                  \begin{pmatrix}A_1&A_2&A_3\\B_1&B_2&B_3\end{pmatrix}.
\]
Take $t$ transcendental over $k$, put $K=\overline{k(t)}$, and define
\[
 E_t:\ v^2=u(u-1)(u-t),\qquad o=\infty,\qquad e=(0,0),
 \qquad X=W_K\times_K E_t.
\]
For every embedding $K\hookrightarrow\mathbb C$, $X_{\mathbb C}$
is a smooth projective fourfold with $\mathrm{CH}_0$ supported on
$\{w_0\}\times E_t$, for any geometric point $w_0$ of $W$.
There is an integral cycle $\gamma\in\mathrm{CH}^2(W)$ such that
\[
 z=\gamma\times([e]-[o])
       \in\mathrm{CH}^3(X_{\mathbb C})_{\mathrm{alg}}
\]
satisfies
\begin{equation}\label{eq:explicit-os-torsion}
             2z=0,\qquad z\ne0,\qquad
             \operatorname{AJ}(z)=0,\qquad z\notin N^3.
\end{equation}
Below we give an exact instruction for choosing $\gamma$.
The displayed variety has no unselected sufficiently-large integer
in its equations.

\begin{proof}
\emph{Integral transfer of the parity obstruction.}
Let $V=\mathbb A^{72}_{\overline{\mathbb Q}}$ be the coefficient
space. A nonempty Zariski open $U\subset V$ carries a smooth
projective family $\pi:\mathcal W\to U$ of these threefolds.
We choose $U$ so that the first-row zero locus in
$\mathbb P^1\times\mathbb P^2$ is finite etale of degree $12$.
Its product with the last $\mathbb P^2$ is then a relative
effective Cartier divisor $\mathcal D\subset\mathcal W$.
The degree follows from
\[
          \int_{\mathbb P^1\times\mathbb P^2}(h+2l)^3=12.
\]
These are open conditions. The general matrices in
\citet[Lemmas 2.3--2.7 and Theorem 3.1, author pp.~4--8]
{cyOttemSuzuki2020} satisfy them. We take $U$ containing one
of their rational matrices $v_0$. Its divisor $D_{v_0}$ is
their $\sum_{j=1}^{12}E_{1,j}$. Their integral congruence (1),
author p.~7, says
\begin{equation}\label{eq:os-parity-all-curves}
 \deg(c/\mathbb P^1)-D_{v_0}\cdot c=0\pmod2
               \quad(c\in\mathrm{CH}_1(W_{v_0})).
\end{equation}
We use this exact conclusion of their degeneration proof as a
known input. Their auxiliary prime is not part of our
definition of $W$.

Let $C$ be an integral curve on $W_{\mathbb C}$. It defines a
point of a relative Hilbert scheme
$\operatorname{Hilb}_P(\mathcal W/U)$ for the fixed Segre
embedding. This scheme is projective over $U$ and defined over
$\overline{\mathbb Q}$. Choose an irreducible component $H$
containing the point. The coefficient tuple is algebraically
independent, so its image in $U$ is the generic point as an
$\overline{\mathbb Q}$-scheme. Thus $H$ dominates $U$.
Its image is closed by properness, so $H$ has a point over $v_0$.
The flat limit defines a one-cycle $c_0$ with integral multiplicities; embedded or
zero-dimensional components do not affect that one-cycle.
The degrees of both
$\mathcal O_{\mathbb P^1}(1)$ and
$\mathcal O_{\mathcal W}(\mathcal D)$ on a flat
one-dimensional family are constant on each connected Hilbert
component. They can be read as the linear coefficients of
Euler characteristics after tensoring by powers of the line bundle.
Consequently
\[
 \deg(C/\mathbb P^1)-D\cdot C
   =\deg(c_0/\mathbb P^1)-D_{v_0}\cdot c_0\equiv0\pmod2.
\]
Additivity proves this for every integral one-cycle on
$W_{\mathbb C}$. No norm, extension degree, or multiplication
of a cycle occurs in this argument.

The same source, author p.~8, gives
$\beta_0\in H_2(W_{v_0},\mathbb Z)$ with
\[
        \deg(\beta_0/\mathbb P^1)=1,\qquad D_{v_0}\cdot\beta_0=0
\]
after a possible sign change. The base $U(\mathbb C)$ is path
connected. Parallel transport and Poincare duality give
$\beta\in H^4(W_{\mathbb C},\mathbb Z)$ with the same two
intersection numbers. Integral cohomology is torsion-free by
Lemma 2.7 or transport from $v_0$. Lemma 2.3 gives
$H^2(W,\mathcal O_W)=0$, so $H^4(W,\mathbb Z)$ consists of
Hodge classes. The parity congruence shows that $\beta$ is
not algebraic.

\emph{An algebraic lift of a class of order two in the defect.}
The subgroup
\[
 B=H^4_{\mathrm{alg}}(W_{\mathbb C},\mathbb Z)
       \subset H^4(W_{\mathbb C},\mathbb Z)
\]
has finite index: the Lefschetz $(1,1)$ theorem makes
$H^2(W,\mathbb Q)$ algebraic, and multiplication by the ample
class of $\mathcal O(1,1,1)|_W$ is an isomorphism
$H^2(W,\mathbb Q)\to H^4(W,\mathbb Q)$.
Let $m>0$ be least with $m\beta\in B$. Parity makes $m$ even.
Choose any algebraic representative
\[
             \operatorname{cl}(\gamma)=m\beta .
\]
We can take $\gamma$ over $k$: for each component of a complex
representative, choose a $k$-point on its Hilbert component for
the fixed variety $W/k$. Algebraic closedness of $k$ gives
geometric connectedness of the component and existence of such
a point. The resulting algebraically equivalent cycle has the
same integral cohomology class.

The element $\operatorname{cl}(\gamma)\otimes\tfrac12$ in
$B\otimes\mathbb Q/\mathbb Z$ is nonzero. Otherwise
$\operatorname{cl}(\gamma)\in2B$; torsion-freeness of $H^4$
would put $(m/2)\beta$ in $B$, contradicting minimality.
Its image in $H^4(W,\mathbb Z)\otimes\mathbb Q/\mathbb Z$ is
zero because $m\beta$ is even. In particular
$\gamma\otimes\tfrac12$ is nonzero in
$\mathrm{CH}^2(W)\otimes\mathbb Q/\mathbb Z$.

\emph{The specified elliptic curve and cycle.}
The Legendre formulas
\[
 \Delta=16t^2(1-t)^2,\qquad
 j(E_t)=256\,\frac{(1-t+t^2)^3}{t^2(1-t)^2}
\]
show smoothness and that $j(E_t)$ is transcendental over $k$.
The class $\tau=[e]-[o]$ is nonzero of order two in
$\operatorname{Pic}^0(E_t)$. Choose a direct summand
$\mathbb Q/\mathbb Z$ of its torsion group taking $1/2$ to $\tau$.
Then $\gamma\otimes\tau\ne0$.
The exact exterior-product injectivity theorem of
\citet[Theorem (0.2), p.~177]{cySchoen2000Exterior}, for
$k\subset\mathbb C$, gives $z=\gamma\times\tau\ne0$.
This cycle has order two and is algebraically trivial.

The Deligne cohomology exact sequence for $W$ has divisible
kernel $J^2(W)$. After tensoring by
$\operatorname{Pic}^0(E_t)_{\mathrm{tors}}\simeq
(\mathbb Q/\mathbb Z)^2$, this kernel contributes zero.
The Deligne class of the product therefore factors through
\[
 H^4_{\mathrm{Hdg}}(W,\mathbb Z)
          \otimes\operatorname{Pic}^0(E_t)_{\mathrm{tors}}.
\]
Its value is $(m\beta)\otimes\tau=0$.
Thus $\operatorname{AJ}(z)=0$. This also verifies the mechanism
in \citet[Proposition 3.1 and proof, author pp.~5--6]
{cySuzuki2020Remark} for our specified field and elliptic curve.

The same chosen $\gamma$ works under every embedding into
$\mathbb C$. Its nonzero tensor class over $k$ is intrinsic.
Under the embedding used to choose it, its Betti class is even,
so comparison makes its mod-two etale cycle class over $k$ zero.
This is an intrinsic class over the algebraically closed field
$k$. Comparison after any other embedding therefore again makes
its integral Betti class even. Schoen's injectivity and the
Deligne-product argument apply unchanged. This proves the
stated quantifier without choosing a new $\gamma$ for each
embedding.

Finally Ottem--Suzuki's Lemma 2.4 gives
$\mathrm{CH}_0(W_{\mathbb C})=\mathbb Z$. In each
$W\times\{e'\}$ we can move every point to $(w_0,e')$, so
$\{w_0\}\times E_t$ supports $\mathrm{CH}_0(X_{\mathbb C})$.
The hypotheses of
\citet[Theorem 1.1, author p.~2]{cySuzuki2020Remark} hold.
That theorem makes the Walker map injective on algebraically
trivial codimension-three torsion. Hence $\psi_W(z)\ne0$;
E.3g gives $N^3\subset\ker\psi_W$, so $z\notin N^3$.

E.3j applies to this same displayed fourfold and proves its
integral family-map converse. We have therefore specified a
function-field counterexample to the literal integral
biconditional in Lead E4. The rational equivalence in E.3c
remains valid. The exact fixed-target Walker-generation
equivalences of E.3o--p remain valid as well; E.3u separately
refutes their universal conditions.
\end{proof}

